%% USPSC-Agradecimentos.tex
\begin{agradecimentos}

	Primeiramente, eu jamais estaria aqui escrevendo estas palavras sem o apoio, cuidado e conhecimento de todos que cruzaram o meu caminho. E honestamente, este espaço é pequeno demais e minhas palavras são imprecisas demais para comunicar o quanto eu tenho a agradecer a todos os seres vivos que fizeram parte desses meus últimos dois anos, mas espero ao menos conseguir transmitir parte da minha gratidão.
	
	Não vejo outra possibilidade senão iniciar os agradecimentos com meus pais. Muito obrigada por serem meu porto seguro em meio ao caos, por sempre me ouvirem e aconselharem a seguir mesmo quando não parecia haver caminhos e por sempre me incentivarem a aprender e crescer. Também não poderia deixar de agradecer toda a alegria, o amor e a torcida da minha família, que cria esse ambiente de pertencimento e acolhimento tão leve e divertido. Gostaria também de dedicar um agradecimento especial às minhas gatas, Morgana e Suri, que, independentemente da distância, sempre melhoram o meu dia. Quero dizer aos meus amigos que a leveza, espontaneidade e aleatoriedades de vocês salvaram muitos dos meus dias. Faço um agradecimento especial aos amigos que fiz em São Carlos, já que minha estadia nessa cidade se tornou muito mais positiva graças à companhia de vocês. Também agradeço de coração ao Lucas pela parceria, incentivo e paciência com meus altos e baixos nessa jornada. Acho essencial que eu agradeça também à minha psicóloga pelo espaço de escuta, suporte e aprendizado.
	
	Além daqueles que me ajudaram a atravessar portas, eu gostaria de agradecer profundamente aqueles que as abriram para mim. Aos meus orientadores, Ivan e Diogo, muito obrigada por me ensinarem tanto sobre física, por me mostrarem como é fazer pesquisa de forma saudável e humana, por criarem espaços acolhedores ao mesmo tempo desafiadores no limite certo para que eu crescesse, e por me apoiarem de todas as formas possíveis. Vocês não só possibilitaram meu mestrado, mas abriram oportunidades para eu transformar em realidade o sonho de ser pesquisadora. Quero também agradecer aos membros do grupo por fazerem da sala 21 um ambiente de trabalho receptivo, divertido e repleto de possibilidades de aprendizado.

    
	
	This study was financed in part by the Coordenação de Aperfeiçoamento de Pessoal de Nível Superior – Brasil (CAPES) – Finance Code 001.
	
\end{agradecimentos}
% ---